CGRAs trade the bit-level flexibility of FPGAs for word-level operations and
fewer configuration bits~\cite{hartenstein,cgra-survey,podobas-survey}.
Classic designs differ mainly in how much storage and communication each
element has: ADRES couples a VLIW processor to a reconfigurable
matrix~\cite{adres}, MorphoSys provides richer communication
resources~\cite{morphosys}, HyCUBE adds reconfigurable multi-hop
paths~\cite{hycube}, PipeRench virtualises pipelined
computations~\cite{piperench}, DySER specialises datapaths inside a
processor~\cite{dyser}, and Plasticine organises compute and memory around
parallel patterns~\cite{plasticine}.

Table~\ref{tab:related} compares the work closest in purpose. CGRA-ME
generates RTL from an architecture description and evaluates CGRA variants
both as standard-cell designs and as FPGA overlays~\cite{cgra-me,cgra-me-ispd}.
QuickDough is a soft CGRA overlay whose PEs hold instruction and data
memories; host data reach on-chip buffers by DMA, and accelerator
performance is estimated from schedule cycles, implementation frequency and
transfer latencies of a Zynq platform~\cite{quickdough}. OpenEdgeCGRA is an
open-hardware CGRA with several registers per PE, a torus interconnect and
per-column program counters, integrated in a microcontroller
platform~\cite{openedgecgra}; Aspros et~al.\ estimate energy and latency of
its kernel executions and compare them with post-synthesis
simulation~\cite{aspros-cgra}. These designs are more capable than ours
(local storage, instruction sequencing, memory access); none of them is
evaluated here, and resources or frequencies obtained with different
devices, precisions, geometries and tool flows are not comparable with ours.
What this paper adds is not a better fabric but a complete accounting of a
deliberately minimal one: every byte, transaction and step on a narrow host
link, and a traceable link between control-logic proofs and timing
exceptions.

\begin{table*}[t]
  \centering
  \caption{Comparison with related CGRA work, as documented in the cited
  primary sources. ``Not reported'' means not found in those sources, not
  absent from the design. Resources and frequencies are deliberately omitted:
  they come from different devices, precisions and flows.}
  \label{tab:related}
  \footnotesize
  \renewcommand{\arraystretch}{1.15}
  \begin{tabularx}{\textwidth}{@{}>{\raggedright\arraybackslash}p{2.65cm}
      *{5}{>{\raggedright\arraybackslash}X}@{}}
    \headtop
    \headrow
    \thead{Work} & \thead{PE state and contexts} & \thead{Communication and memory}
      & \thead{Host interface} & \thead{Mapping} & \thead{Validation reported} \\
    \headbottom
    CGRA-ME~\cite{cgra-me,cgra-me-ispd}
      & Described per architecture (XML)
      & Architecture-defined
      & Not reported
      & Integrated mapper
      & Standard-cell and FPGA-overlay implementation \\
    QuickDough~\cite{quickdough}
      & Instruction memory and multi-port data memory per PE
      & Torus; on-chip I/O buffers
      & DMA from an ARM host (Zynq)
      & Compile-time scheduling of loop DFGs
      & Analytical performance from schedule, $f_{\max}$ and transfer latency \\
    OpenEdgeCGRA~\cite{openedgecgra}
      & Four registers plus output register; per-column program counter
      & Torus; column memory ports
      & Microcontroller integration (X-HEEP)
      & Kernel assembly, behavioural simulator
      & Not reported in the repository; post-synthesis comparison in~\cite{aspros-cgra} \\
    This work
      & One result register; one configuration
      & Mesh with north/west edge injection; optional two $1024\times16$ banks
      & UART, 115\,200\,Bd
      & Fixed library mappings with host reduction
      & Exact counters, RTL replay, formal safety proofs, STA; no board \\
    \bottomrule
  \end{tabularx}
\end{table*}

